% Ad Familiares 14.14 --- headnote
% ad-familiares-14-14, 25 January 49 BC, Minturnae

\textit{Cicero to his wife Terentia and daughter Tullia,
written from Minturnae on the eighth day before the
Kalends of February 49 BC (the manuscript dateline
mis-records the month as \textit{Quint.}; it is plainly
February, the day after the Formiae note). Cicero is
still on his way south to take up the coastal command,
travelling with his son Marcus and his brother Quintus
(and with the younger Quintus and the friend Rufus),
who add their greetings at the close.}

\textit{The letter follows on directly from
\textit{Fam.}~14.18 of the previous day: the same
question, sharper. The two alternatives have hardened
into one fork --- if Caesar enters Rome with restraint
the household can stay; if not, even Dolabella will be
no protection. The fear now added is of being cut off
on the road, of famine in the city, of the women of
their standing having already gone. The good news at
the end --- Labienus has come over to the Pompeian
side; Piso (the father-in-law of Caesar himself) is
leaving the city and calling his son-in-law's act a
crime --- is the running political bulletin the family
is being kept on. The household consultation list ---
Pomponius, Camillus --- is the same one the Atticus
letters use.}
